% ============================================================================
% MUSTERLÖSUNG DS 8 - Gruppe A + B: Mi persona favorita
% ============================================================================

\documentclass[11pt,a4paper]{article}

\usepackage[utf8]{inputenc}
\usepackage[T1]{fontenc}
\usepackage[ngerman]{babel}
\usepackage[left=2cm,right=2cm,top=1.8cm,bottom=1.5cm]{geometry}
\usepackage{helvet}
\usepackage{xcolor}
\usepackage{tabularx}
\usepackage{array}
\usepackage{enumitem}
\usepackage{tcolorbox}
\usepackage{fancyhdr}
\usepackage{lastpage}

\renewcommand{\familydefault}{\sfdefault}

\definecolor{spanisch}{HTML}{c41e3a}
\definecolor{grau}{HTML}{666666}
\definecolor{hellgrau}{HTML}{f5f5f5}
\definecolor{gruppeB}{HTML}{1565c0}
\definecolor{loesung}{HTML}{c62828}

\pagestyle{fancy}
\fancyhf{}
\fancyhead[L]{\small\textcolor{grau}{Spanisch Spätbeginner -- MUSTERLÖSUNG}}
\fancyhead[R]{\small\textcolor{grau}{ML-08}}
\fancyfoot[C]{\small\thepage/\pageref{LastPage}}
\renewcommand{\headrulewidth}{0.5pt}

\newcommand{\lsg}[1]{\textcolor{loesung}{\textbf{#1}}}
\newcommand{\es}[1]{\textit{#1}}

\newtcolorbox{gruppenbox}[1][]{
    colback=white,
    colframe=spanisch,
    fonttitle=\bfseries\color{white},
    colbacktitle=spanisch,
    title=#1,
    boxrule=1pt,
    arc=0pt,
    left=5pt, right=5pt, top=5pt, bottom=5pt
}

\newtcolorbox{gruppeBbox}[1][]{
    colback=white,
    colframe=gruppeB,
    fonttitle=\bfseries\color{white},
    colbacktitle=gruppeB,
    title=#1,
    boxrule=1pt,
    arc=0pt,
    left=5pt, right=5pt, top=5pt, bottom=5pt
}

\begin{document}

% ============================================================================
% KOPFBEREICH
% ============================================================================
\begin{center}
{\Large\textbf{\textcolor{loesung}{MUSTERLÖSUNG}}}\\[0.3em]
{\large DS 8 -- Mi persona favorita (Creatividad y Juego)}
\end{center}

\vspace{0.5em}
\hrule
\vspace{1em}

% ============================================================================
% GRUPPE A
% ============================================================================
\begin{gruppenbox}[Gruppe A: Mi persona favorita (32 Punkte)]

\textbf{Kasten 1: Adjetivos para describir personas}

\begin{tabularx}{\textwidth}{@{}|X|X|X|@{}}
\hline
\textbf{Positiv} & \textbf{Neutral} & \textbf{Charakter} \\
\hline
simpático/-a & tranquilo/-a & trabajador/-a \\
inteligente & serio/-a & responsable \\
divertido/-a & reservado/-a & organizado/-a \\
amable & \lsg{callado/-a (still)} & creativo/-a \\
paciente & \lsg{tímido/-a (schüchtern)} & \lsg{honesto/-a (ehrlich)} \\
\hline
\end{tabularx}

\vspace{0.8em}

\textbf{Aufgabe 1: Wortschatz-Staffel (6 P.)}

Mögliche Adjektive (je 0,5 P., max. 6 P.):
\begin{itemize}[leftmargin=1.5em, nosep]
    \item \lsg{simpático, inteligente, divertido, amable, paciente, creativo, trabajador, responsable, organizado, tranquilo, serio, reservado, cariñoso, generoso, tolerante...}
\end{itemize}

\vspace{0.5em}

\textbf{Aufgabe 2: ser ergänzen (6 P.)}
\begin{enumerate}[label=\alph*), leftmargin=2em, nosep]
    \item Mi padre \lsg{es} muy trabajador.
    \item Mis abuelos \lsg{son} simpáticos y pacientes.
    \item ¿Cómo \lsg{es} tu mejor amigo?
    \item Nosotros \lsg{somos} estudiantes creativos.
    \item Ella \lsg{es} inteligente y responsable.
    \item Yo \lsg{soy} una persona organizada.
\end{enumerate}

\vspace{0.5em}

\textbf{Aufgabe 3: Konnektoren (8 P.)}
\begin{enumerate}[label=\alph*), leftmargin=2em, nosep]
    \item \lsg{Mi hermana es simpática y también es muy inteligente.}
    \item \lsg{Mi profesor es serio, pero es muy amable.}
    \item \lsg{Me gusta mi abuela porque es muy paciente.}
    \item \lsg{Mi amigo es divertido y además es creativo.}
\end{enumerate}

\vspace{0.5em}

\textbf{Aufgabe 4: Poster-Vorbereitung (8 P.) -- Beispiellösung:}

\lsg{Mi persona favorita es mi madre. Es mi madre y tiene 48 años.}

\lsg{Es muy amable y también muy paciente. Además, es una persona muy trabajadora y organizada.}

\lsg{Me gusta mucho porque siempre me ayuda con mis problemas. Siempre tiene tiempo para mí y me escucha.}

\vspace{0.5em}

\textbf{Aufgabe 5: Gallery Walk (4 P.)}

Bewertung nach Kriterien: Sprache, Kreativität, Gesamteindruck (je 1 P. pro Poster = max. 3 Poster bewertet)

\end{gruppenbox}

\vspace{1em}

% ============================================================================
% GRUPPE B
% ============================================================================
\begin{gruppeBbox}[Gruppe B: Mi persona favorita -- Nivel avanzado (35 Punkte)]

\textbf{Kasten 1: Adjetivos (erweitert)}

\begin{tabularx}{\textwidth}{@{}|X|X|X|@{}}
\hline
\textbf{Positiv} & \textbf{Beruflich} & \textbf{Persönlichkeit} \\
\hline
simpático/-a, amable & trabajador/-a & introvertido/-a \\
inteligente, listo/-a & responsable & extrovertido/-a \\
divertido/-a, gracioso/-a & organizado/-a & optimista \\
paciente, tolerante & profesional & pesimista \\
cariñoso/-a & \lsg{dedicado/-a} & ambicioso/-a \\
generoso/-a & \lsg{eficiente} & \lsg{sensible} \\
\hline
\end{tabularx}

\vspace{0.8em}

\textbf{Aufgabe 1: Wortschatz nach Kategorien (6 P.)}

\begin{tabularx}{\textwidth}{@{}|X|X|X|@{}}
\hline
\textbf{Positiv} & \textbf{Neutral/Negativ} & \textbf{Beruf/Charakter} \\
\hline
\lsg{simpático} & \lsg{serio} & \lsg{trabajador} \\
\lsg{amable} & \lsg{tímido} & \lsg{responsable} \\
\lsg{divertido} & \lsg{reservado} & \lsg{organizado} \\
\lsg{paciente} & \lsg{pesimista} & \lsg{ambicioso} \\
\hline
\end{tabularx}

\vspace{0.5em}

\textbf{Aufgabe 2: ser / estar / parecer (6 P.)}
\begin{enumerate}[label=\alph*), leftmargin=2em, nosep]
    \item Mi jefe \lsg{parece} muy estricto, pero en realidad \lsg{es} muy amable.
    \item Hoy mi madre \lsg{está} cansada porque ha trabajado mucho.
    \item Mis amigos \lsg{son} personas muy divertidas.
    \item Él \lsg{parece} serio, pero cuando lo conoces, \lsg{es} simpático.
    \item ¿Cómo \lsg{es} tu profesor de español?
    \item Ella siempre \lsg{está} de buen humor.
\end{enumerate}

\vspace{0.5em}

\textbf{Aufgabe 3: Komplexe Sätze (8 P.)}
\begin{enumerate}[label=\alph*), leftmargin=2em, nosep]
    \item \lsg{Mi madre parece estricta, pero es muy paciente.} \\
          (Alternativ: \lsg{Aunque mi madre parece estricta, es muy paciente.})
    \item \lsg{Mi amigo es muy inteligente, por eso siempre me ayuda.}
    \item \lsg{Aunque mi padre trabaja mucho, siempre tiene tiempo para mí.}
    \item \lsg{Ya que mi abuela es muy sabia, le pido consejo a menudo.}
\end{enumerate}

\vspace{0.5em}

\textbf{Aufgabe 4: Poster-Text (10 P.) -- Beispiellösung:}

\lsg{Mi persona favorita es mi abuelo. Se llama Antonio y tiene 72 años. Es una persona muy especial para mí.}

\lsg{Mi abuelo es extremadamente paciente y bastante sabio. Además, es muy cariñoso y siempre está de buen humor. Aunque parece serio, en realidad es muy divertido.}

\lsg{Le gusta contar historias de su juventud y también le gusta cocinar. Siempre prepara mi comida favorita cuando lo visito.}

\lsg{Es importante para mí porque siempre me escucha y me da buenos consejos. Por ejemplo, cuando tenía problemas en el colegio, él me ayudó a encontrar una solución.}

\vspace{0.5em}

\textbf{Aufgabe 5: Escape-Room-Rätsel (5 P.)}

\textbf{Rätsel 1:}
\begin{itemize}[leftmargin=1.5em, nosep]
    \item Adjektive: \lsg{Simpático, Inteligente, Creativo, Paciente, Responsable, Divertido}
    \item Anfangsbuchstaben: \lsg{S -- I -- C -- P -- R -- D}
    \item 3. Buchstabe = C = \lsg{3} $\rightarrow$ \textbf{Code-Ziffer 1: \lsg{3}}
\end{itemize}

\textbf{Rätsel 2:}
\begin{itemize}[leftmargin=1.5em, nosep]
    \item \es{ser}: \lsg{es} (1), \lsg{somos} (2) = \textbf{2 Mal}
    \item \es{estar}: estoy, está, estás = 3 Mal (zählt nicht!)
    \item \textbf{Code-Ziffer 2: \lsg{2}}
\end{itemize}

\textbf{Rätsel 3:}
\begin{itemize}[leftmargin=1.5em, nosep]
    \item A (médico) = \lsg{3} (paciente)
    \item B (profesor) = \lsg{4} (organizado)
    \item C (artista) = \lsg{1} (creativo)
    \item D (bombero) = \lsg{2} (valiente)
    \item A + B = 3 + 4 = \textbf{Code-Ziffer 3: \lsg{7}}
\end{itemize}

\textbf{Gesamtcode: \lsg{3 -- 2 -- 7}}

\end{gruppeBbox}

\vspace{1em}
\hrule
\vspace{0.3em}

\textbf{Bewertungshinweise:}
\begin{itemize}[leftmargin=1.5em, nosep]
    \item Gruppe A: 32 Punkte gesamt
    \item Gruppe B: 35 Punkte gesamt
    \item Bei Poster-Texten: Inhalt und Struktur wichtiger als perfekte Grammatik
    \item Konnektoren: Auch alternative korrekte Verbindungen akzeptieren
    \item Halbe Punkte möglich bei teilweise richtigen Antworten
    \item Wortschatz-Staffel: Auch nicht aufgelistete korrekte Adjektive akzeptieren
\end{itemize}

\vspace{0.5em}
\textbf{Differenzierungshinweise:}
\begin{itemize}[leftmargin=1.5em, nosep]
    \item \textbf{Gruppe A:} Fokus auf Grundstrukturen (\es{ser} + Adj., einfache Konnektoren)
    \item \textbf{Gruppe B:} Erweiterung um \es{parecer}, \es{estar}, komplexe Konnektoren (\es{aunque}, \es{ya que})
    \item \textbf{Escape-Room:} Identisch für beide Gruppen (gemeinsames Erlebnis)
\end{itemize}

\end{document}
